% Unite: papier 30, introduction; traduction francaise depuis la source allemande auditee.
\setcounter{footnote}{0}

\section*{30. Construction abstraite de la théorie des idéaux dans les corps de nombres et de fonctions algébriques}
\noindent\emph{Math. Ann. 96 (1927), p. 26--61.}

\begin{center}
Par\\[0.5em]
\textbf{Emmy Noether à Göttingen.}
\end{center}

Dans ce qui suit on donne une caractérisation abstraite de tous les anneaux dont la théorie des idéaux coïncide avec la théorie des idéaux de toutes les grandeurs entières du corps de nombres algébriques -- c'est-à-dire dont les idéaux se représentent de façon univoque comme produits de puissances d'idéaux premiers. À ces anneaux appartiennent encore, comme exemples les plus connus, l'anneau de toutes les grandeurs entières d'un corps de fonctions algébriques d'une indéterminée -- plus généralement de plusieurs indéterminées, dès qu'avec Kronecker on se restreint aux idéaux de dimension maximale, donc que l'on passe à une certaine localisation, le domaine fonctionnel.

Les axiomes qui sont complètement équivalents à cette théorie des idéaux sont, pour l'anneau commutatif pris pour base, les suivants:\footnote{Pour la définition des notions fondamentales de la théorie des idéaux, comparer par exemple E. Noether, \emph{Idealtheorie in Ringbereichen}, Math. Ann. 83 (1921), p. 24--66 (cité: \emph{Idealtheorie}). D'ailleurs les notions les plus essentielles sont aussi formulées brièvement dans le présent travail, surtout aux \S{}\S{}1, 4 et 5. Pour être complète, la présentation contient aussi des éléments qui ne sont pas nécessaires aux fins immédiates de ce travail.}

\begin{description}
\item[I. Condition de chaîne des diviseurs.] Toute chaîne d'idéaux dans laquelle chaque idéal est un diviseur propre du précédent s'interrompt au bout d'un nombre fini de pas -- autrement dit: pour toute chaîne de diviseurs d'idéaux, il existe un indice à partir duquel tous les idéaux deviennent égaux.
\item[II. Condition de chaîne des multiples modulo tout idéal différent de l'idéal nul.] Toute chaîne d'idéaux -- qui sont tous diviseurs d'un idéal fixe différent de l'idéal nul -- dans laquelle chaque idéal est un multiple propre du précédent s'interrompt au bout d'un nombre fini de pas.
\item[III.] Existence de l'élément unité de la multiplication.
\item[IV.] Anneau sans diviseurs de zéro.
\item[V. Clôture intégrale dans le corps des fractions.] Tout élément du corps des fractions qui est entier relativement à l'anneau appartient à l'anneau.
\end{description}

À partir de ces axiomes je développe pas à pas une théorie des idéaux de plus en plus restreinte, jusqu'à celle que l'on cherche; et je montre réciproquement que tous les axiomes y sont effectivement satisfaits.

De la condition de chaîne des diviseurs I il résulte, comme je l'ai montré (\emph{Idealtheorie} \S{}4), la représentation d'un idéal comme plus petit commun multiple d'un nombre fini d'idéaux primaires appartenant à des idéaux premiers différents. Je répète ici brièvement cette démonstration (\S{}6), parce qu'elle se simplifie en utilisant la notion de quotient d'idéaux, et parce que je veux corriger une inadvertance: la démonstration originelle utilise en effet, en un point, l'existence non supposée de l'élément unité de la multiplication.\footnote{P. Urysohn a attiré mon attention sur cette inadvertance; ma modification de la démonstration était toutefois un peu plus compliquée que celle communiquée ici, qui est due à E. Artin. Celui-ci avait déjà comblé plus tôt la lacune pour lui-même et avait aussi, indépendamment de moi, envisagé la simplification au moyen du quotient d'idéaux. -- Une autre simplification, qui évite l'introduction de la ``représentation réduite'', je la tire d'une question apparentée chez W. Krull: \emph{Algebraische Theorie der Ringe III}, Math. Ann. 92 (1924), p. 181--213, théorème I.}

L'hypothèse de la condition de chaîne des multiples a pour effet (\S{}7) que, dans l'anneau quotient par tout idéal différent de l'idéal nul, un idéal premier ne peut posséder aucun diviseur propre, à l'exception de l'idéal unité comprenant tous les éléments. Les composantes primaires de ces idéaux sont donc déterminées de façon univoque, à l'exception de celles qui appartiennent à l'idéal unité. Sous une forme un peu moins précise, ce résultat se trouve déjà chez Masazo Sono;\footnote{Masazo Sono, \emph{On the Reduction of Ideals}, Memoirs of the College of Science, Kyoto University, Series A, 7 (1924), p. 191--204.} son hypothèse de l'existence d'une série de composition dans l'anneau quotient est équivalente à la ``condition de double chaîne'' dans cet anneau, comme je le montre brièvement à la fin (\S{}10). Par condition de double chaîne j'entends ici la validité des conditions de chaîne des diviseurs et des multiples sans restriction aux idéaux différents de l'idéal nul.

Si l'axiome III -- existence de l'élément unité -- est encore rempli, l'idéal unité ne se présente pas comme idéal premier associé; les composantes primaires sont donc déterminées de façon univoque; elles deviennent étrangères deux à deux et leur plus petit commun multiple est égal au produit. Les axiomes I, II, III sont remplis pour tous les ordres finis d'un corps de nombres algébriques; ici Dedekind avait déjà énoncé (Dirichlet--Dedekind, \emph{Zahlentheorie}, 3e éd., 11e suppl., \S{}172), sans démonstration entièrement développée, la représentation univoque d'un idéal comme produit d'idéaux d'une seule espèce deux à deux étrangers (composantes primaires).

De l'hypothèse IV -- absence de diviseurs de zéro -- il résulte que les différentes puissances d'un idéal différent de l'idéal nul et de l'idéal unité sont toutes différentes, et que le corps des fractions existe. Si enfin l'axiome V de la clôture intégrale dans le corps des fractions est encore rempli, alors les idéaux primaires -- déterminés univoquement à cause de IV -- deviennent des puissances d'idéaux premiers, ce qui donne le théorème usuel de décomposition (\S{}8). Cette dernière démonstration repose sur une conséquence de la clôture intégrale tirée déjà par Dedekind (3e éd., \S{}172) dans le cas du corps de nombres: on peut représenter un élément véritablement fractionnaire comme quotient d'éléments entiers de telle sorte que même le carré du numérateur ne soit pas divisible par le dénominateur. Dans les présentations ultérieures, cette conséquence est remplacée -- également comme conséquence de la clôture intégrale -- par le théorème de Gauss généralisé, beaucoup moins transparent, ou par des théorèmes de modules correspondants, un peu plus forts; tous ces théorèmes supposent, pour leur application, des éléments de base, tandis qu'ici on travaille avec les idéaux eux-mêmes.

Dans la réciproque \hbox{\(\S 9\)}, les axiomes autres que l'axiome V s'obtiennent presque directement; celui-ci s'obtient en démontrant qu'un élément véritablement fractionnaire peut toujours se représenter de telle sorte qu'aucune puissance du numérateur ne soit divisible par le dénominateur, ce qui est donc encore un renforcement de la conséquence primitivement tirée de V. Mais cette démonstration ne réussit qu'en tirant de la représentation des idéaux les conclusions usuelles: représentation par idéaux principaux modulo un idéal fixe et théorie des idéaux fractionnaires, c'est-à-dire des quotients d'idéaux dans le corps. Le fait que la clôture intégrale résulte de la représentation par puissances d'idéaux premiers était déjà connu de Dedekind, comme le montre une remarque (3e éd., \S{}172, p. 522 en bas).

Les axiomes I à V impliquent que les quatre décompositions en général distinctes -- en plus petits idéaux étrangers deux à deux, en idéaux mutuellement premiers, en plus grandes composantes primaires et en idéaux irréductibles -- coïncident; et réciproquement, de cette coïncidence et des axiomes I--IV résulte la clôture intégrale. Pour les anneaux avec diviseurs de zéro, le fait que les autres axiomes soient remplis -- V étant à définir comme clôture intégrale dans l'anneau total des fractions -- n'implique pas nécessairement la coïncidence des décompositions, comme le montre l'exemple de l'anneau quotient par certains idéaux d'un ordre fini \hbox{\(\S 9\)}.

Au \S{}1 j'esquisse la théorie des grandeurs entières relativement à un anneau, jusqu'à la conséquence déduite par Dedekind de la clôture intégrale. Au \S{}2 je montre comment les conditions de chaîne se transfèrent de l'anneau aux modules formés de formes linéaires, avec cet anneau comme domaine de coefficients et de multiplicateurs.\footnote{Cette forme des théorèmes de modules provient de Prüfer (Neue Begründung der algebraischen Zahlentheorie, \S{}2, Math. Ann. 94 (1924), p. 198--243) et est plus transparente que le transfert usuel des théorèmes de base.} J'en tire (\S{}3) la conséquence que, lors du passage aux grandeurs entières d'un corps d'extension algébrique (de première espèce), les axiomes I à IV se conservent pour tout ordre, tandis que V vaut aussi pour l'ordre formé de toutes les grandeurs entières. Ce fait -- naturellement connu pour le corps de nombres algébriques -- permet de subordonner les corps de fonctions mentionnés au début; en particulier, par la simple démonstration que le domaine fonctionnel dérivé des polynômes entiers en plusieurs indéterminées satisfait aux axiomes, la théorie des idéaux de Kronecker est elle aussi pleinement intégrée.\footnote{La différence n'apparaît donc qu'avec les théorèmes sur les discriminants, du fait que l'anneau quotient du domaine fonctionnel par un nombre premier devient un corps imparfait, et que des extensions de seconde espèce peuvent ainsi se présenter.}

La théorie des idéaux développée ensuite ne présuppose pas ces \S{}\S{}2 et 3, qui ne servent qu'à l'intégration. Aux \S{}4 et \S{}5 sont données les bases indépendantes des axiomes I à V; ainsi ressort plus nettement que dans la fondation originelle quelles parties de la théorie sont indépendantes d'hypothèses de finitude.
